\documentclass[10pt]{article}
\usepackage[utf8]{inputenc}
\usepackage[vietnamese]{babel}
\usepackage{amsmath,amssymb,amsfonts}
\usepackage{geometry}
\geometry{
    paperwidth=8.5in,
    paperheight=11in,
    top=0.6in,
    bottom=0.6in,
    left=0.65in,
    right=0.65in
}
\usepackage{xcolor}
\definecolor{stewartcyan}{RGB}{0, 130, 185}
\definecolor{stewartgray}{gray}{0.4}

\usepackage{fancyhdr}
\usepackage{tikz}
\usepackage{tabularx}
\usepackage{microtype}

\pagestyle{fancy}
\fancyhf{}
\renewcommand{\headrulewidth}{0pt}
\fancyhead[L]{\textbf{\textsf{224}}\qquad \textbf{\textsf{CHƯƠNG 3}}\quad \textsf{Các quy tắc tính đạo hàm}}
\fancyfoot[C]{\fontsize{6.5}{8}\selectfont\color{stewartgray}Copyright 2021 Cengage Learning. All Rights Reserved. May not be copied, scanned, or duplicated, in whole or in part. WCN 02-200-203\\[1pt]Editorial review has deemed that any suppressed content does not materially affect the overall learning experience. Cengage Learning reserves the right to remove additional content at any time if subsequent rights restrictions require it.}

\newcommand{\casicon}{%
  \tikz[baseline=-0.7ex]{%
    \draw[stewartcyan, fill=cyan!8, rounded corners=0.8pt, line width=0.6pt] (0,0) rectangle (9.5pt, 8pt);
    \draw[stewartcyan, line width=0.6pt] (1.8pt, 2pt) .. controls (3.8pt, 7pt) and (5.8pt, 1pt) .. (7.8pt, 6pt);
  }%
}

\newcommand{\coldivider}{%
  \par\vspace{0.45em}
  {\color{stewartcyan}\hrule height 0.6pt}
  \par\vspace{0.45em}
}

\setlength{\parindent}{0pt}
\setlength{\parskip}{0pt}

\begin{document}

% Tiêu đề mục Bài tập chuẩn Stewart (4 vạch kẻ hai bên)
\noindent
\begin{tikzpicture}
  \node[anchor=west, inner sep=0] (title) at (48pt, 0) {%
    {\fontsize{16}{18}\selectfont\bfseries\color{stewartcyan}3.6}\hspace{8pt}%
    {\color{stewartcyan}\vrule width 0.8pt height 13pt depth 3pt}\hspace{8pt}%
    {\fontsize{16}{18}\selectfont\bfseries\textsf{Bài tập}}%
  };
  \foreach \y in {-4.5, -1.5, 1.5, 4.5} {
    \draw[stewartcyan, line width=0.5pt] (0pt, \y) -- (40pt, \y);
    \draw[stewartcyan, line width=0.5pt] ([xshift=8pt]title.east |- 0, \y) -- (\linewidth, \y);
  }
\end{tikzpicture}
\vspace{0.8em}

\noindent
\begin{minipage}[t]{0.485\textwidth}
\textbf{1.} Giải thích tại sao hàm số logarit tự nhiên $y = \ln x$ được sử dụng thường xuyên hơn nhiều trong giải tích so với các hàm số logarit khác $y = \log_b x$.

\vspace{0.5em}
\noindent
{\bfseries\color{stewartcyan}2--26}\enspace \textbf{Tính đạo hàm của hàm số.}

\vspace{0.35em}
\noindent
\textbf{2.} $f(t) = \ln(3 + t^2)$

\vspace{0.25em}
\noindent
\begin{tabularx}{\linewidth}{@{}X@{\hspace{0.5em}}X@{}}
\textbf{3.} $f(x) = \ln(x^2 + 3x + 5)$ & \textbf{4.} $f(x) = x\ln x - x$ \\[0.35em]
\textbf{5.} $f(x) = \sin(\ln x)$ & \textbf{6.} $f(x) = \ln(\sin^2 x)$ \\[0.35em]
\textbf{7.} $f(x) = \ln\dfrac{1}{x}$ & \textbf{8.} $y = \dfrac{1}{\ln x}$ \\[0.6em]
\textbf{9.} $g(x) = \ln(xe^{-2x})$ & \textbf{10.} $g(t) = \sqrt{1 + \ln t}$ \\[0.45em]
\textbf{11.} $F(t) = (\ln t)^2 \sin t$ & \textbf{12.} $p(t) = \ln\sqrt{t^2 + 1}$ \\[0.45em]
\textbf{13.} $y = \log_8(x^2 + 3x)$ & \textbf{14.} $y = \log_{10}\sec x$ \\[0.45em]
\textbf{15.} $F(s) = \ln\ln s$ & \textbf{16.} $P(v) = \dfrac{\ln v}{1 - v}$ \\[0.65em]
\textbf{17.} $T(z) = 2^z \log_2 z$ & \textbf{18.} $g(t) = \ln\dfrac{t(t^2 + 1)^4}{\sqrt[3]{2t - 1}}$ \\[0.85em]
\textbf{19.} $y = \ln|3 - 2x^5|$ & \textbf{20.} $y = \ln(\csc x - \cot x)$ \\[0.35em]
\textbf{21.} $y = \ln(e^{-x} + xe^{-x})$ & \textbf{22.} $g(x) = e^{x^2 \ln x}$ \\[0.35em]
\textbf{23.} $h(x) = e^{x^2 + \ln x}$ & \textbf{24.} $y = \ln\sqrt{\dfrac{1 + 2x}{1 - 2x}}$ \\[0.85em]
\textbf{25.} $y = \ln\dfrac{x^a}{b^x}$ & \textbf{26.} $y = \log_2(x\log_5 x)$
\end{tabularx}

\coldivider

\noindent
\textbf{27.} Chứng minh rằng $\dfrac{d}{dx}\ln(x + \sqrt{x^2 + 1}) = \dfrac{1}{\sqrt{x^2 + 1}}$.

\vspace{0.45em}
\noindent
\textbf{28.} Chứng minh rằng $\dfrac{d}{dx}\ln\sqrt{\dfrac{1 - \cos x}{1 + \cos x}} = \csc x$.

\coldivider

\noindent
{\bfseries\color{stewartcyan}29--32}\enspace \textbf{Tìm $y'$ và $y''$.}

\vspace{0.35em}
\noindent
\begin{tabularx}{\linewidth}{@{}X@{\hspace{0.5em}}X@{}}
\textbf{29.} $y = \sqrt{x}\ln x$ & \textbf{30.} $y = \dfrac{\ln x}{1 + \ln x}$ \\[0.6em]
\textbf{31.} $y = \ln|\sec x|$ & \textbf{32.} $y = \ln(1 + \ln x)$
\end{tabularx}

\coldivider

\noindent
{\bfseries\color{stewartcyan}33--36}\enspace \textbf{Tính đạo hàm của $f$ và tìm tập xác định của $f$.}

\vspace{0.35em}
\noindent
\begin{tabularx}{\linewidth}{@{}X@{\hspace{0.5em}}X@{}}
\textbf{33.} $f(x) = \dfrac{x}{1 - \ln(x - 1)}$ & \textbf{34.} $f(x) = \sqrt{2 + \ln x}$ \\[0.7em]
\textbf{35.} $f(x) = \ln(x^2 - 2x)$ & \textbf{36.} $f(x) = \ln\ln\ln x$
\end{tabularx}

\vspace{0.6em}
\noindent
\textbf{37.} Nếu $f(x) = \ln(x + \ln x)$, hãy tìm $f'(1)$.

\vspace{0.35em}
\noindent
\textbf{38.} Nếu $f(x) = \cos(\ln x^2)$, hãy tìm $f'(1)$.

\end{minipage}%
\hfill
\begin{minipage}[t]{0.485\textwidth}
\noindent
{\bfseries\color{stewartcyan}39--40}\enspace \textbf{Tìm phương trình tiếp tuyến của đường cong tại điểm đã cho.}

\vspace{0.35em}
\noindent
\textbf{39.} $y = \ln(x^2 - 3x + 1), \quad (3, 0)$

\vspace{0.3em}
\noindent
\textbf{40.} $y = x^2\ln x, \quad (1, 0)$

\coldivider

\noindent
\casicon\enspace \textbf{41.} Nếu $f(x) = \sin x + \ln x$, hãy tìm $f'(x)$. Kiểm tra xem câu trả lời của bạn có hợp lý không bằng cách so sánh đồ thị của $f$ và $f'$.

\vspace{0.4em}
\noindent
\casicon\enspace \textbf{42.} Tìm các phương trình tiếp tuyến của đường cong $y = (\ln x)/x$ tại các điểm $(1, 0)$ và $(e, 1/e)$. Minh họa bằng cách vẽ đồ thị của đường cong và các tiếp tuyến đó.

\vspace{0.4em}
\noindent
\textbf{43.} Cho $f(x) = cx + \ln(\cos x)$. Với giá trị nào của $c$ thì $f'(\pi/4) = 6$?

\vspace{0.4em}
\noindent
\textbf{44.} Cho $f(x) = \log_b(3x^2 - 2)$. Với giá trị nào của $b$ thì $f'(1) = 3$?

\coldivider

\noindent
{\bfseries\color{stewartcyan}45--56}\enspace \textbf{Sử dụng đạo hàm logarit để tìm đạo hàm của hàm số.}

\vspace{0.35em}
\noindent
\begin{tabularx}{\linewidth}{@{}X@{\hspace{0.5em}}X@{}}
\textbf{45.} $y = (x^2 + 2)^2(x^4 + 4)^4$ & \textbf{46.} $y = \dfrac{e^{-x}\cos^2 x}{x^2 + x + 1}$ \\[0.75em]
\textbf{47.} $y = \sqrt{\dfrac{x - 1}{x^4 + 1}}$ & \textbf{48.} $y = \sqrt{x}\,e^{x^2 - 2x}(x + 1)^{2/3}$ \\[0.75em]
\textbf{49.} $y = x^x$ & \textbf{50.} $y = x^{1/x}$ \\[0.45em]
\textbf{51.} $y = x^{\sin x}$ & \textbf{52.} $y = (\sqrt{x})^x$ \\[0.45em]
\textbf{53.} $y = (\cos x)^x$ & \textbf{54.} $y = (\sin x)^{\ln x}$ \\[0.45em]
\textbf{55.} $y = x^{\ln x}$ & \textbf{56.} $y = (\ln x)^{\cos x}$
\end{tabularx}

\coldivider

\noindent
\textbf{57.} Tìm $y'$ nếu $y = \ln(x^2 + y^2)$.

\vspace{0.25em}
\noindent
\textbf{58.} Tìm $y'$ nếu $x^y = y^x$.

\vspace{0.25em}
\noindent
\textbf{59.} Tìm công thức cho $f^{(n)}(x)$ nếu $f(x) = \ln(x - 1)$.

\vspace{0.25em}
\noindent
\textbf{60.} Tìm $\dfrac{d^9}{dx^9}(x^8\ln x)$.

\vspace{0.25em}
\noindent
\textbf{61.} Sử dụng định nghĩa đạo hàm để chứng minh rằng
\[
\lim_{x \to 0}\frac{\ln(1 + x)}{x} = 1
\]

\vspace{-0.2em}
\noindent
\textbf{62.} Chứng minh rằng $\lim\limits_{n \to \infty}\left(1 + \dfrac{x}{n}\right)^n = e^x$ với mọi $x > 0$.

\coldivider

\noindent
{\bfseries\color{stewartcyan}63--78}\enspace \textbf{Tính đạo hàm của hàm số. Rút gọn nếu có thể.}

\vspace{0.35em}
\noindent
\begin{tabularx}{\linewidth}{@{}X@{\hspace{0.5em}}X@{}}
\textbf{63.} $f(x) = \sin^{-1}(5x)$ & \textbf{64.} $g(x) = \sec^{-1}(e^x)$ \\[0.45em]
\textbf{65.} $y = \tan^{-1}\sqrt{x - 1}$ & \textbf{66.} $y = \tan^{-1}(x^2)$ \\[0.45em]
\textbf{67.} $y = (\tan^{-1}x)^2$ & \textbf{68.} $g(x) = \arccos\sqrt{x}$ \\[0.45em]
\textbf{69.} $h(x) = (\arcsin x)\ln x$ & \textbf{70.} $g(t) = \ln(\arctan(t^4))$ \\[0.45em]
\textbf{71.} $f(z) = e^{\arcsin(z^2)}$ & \textbf{72.} $y = \tan^{-1}(x - \sqrt{1 + x^2})$
\end{tabularx}

\end{minipage}

\end{document}